% =====================================================================
% Section: Conclusion  (outline.md §5)
% Draft v1 (2026-06-27). Summarizes contributions and results, notes the
%   transferability of DPR as a generic content-routing module, and future work.
% Consistent with the honest framing of the body (matches/slightly surpasses
%   CATANet at comparable compute; ablation is mechanism-level under shared init).
% No numeric claims beyond those established in Sec. 3.
% =====================================================================
\section{Conclusion}
\label{sec:conclusion}

We presented Dynamic Prototype Routing (DPR), a drop-in replacement for the
center-based Token Aggregation Block of CATANet that keeps the backbone, the tensor
interface, and the lightweight budget unchanged. DPR generates semantic prototypes
directly from the current tokens, decouples their confirmation through a learnable
query bank, matches tokens back with a learnable temperature, and propagates an
explicit per-token confidence through ordering, gating, and a soft fallback, with a
usage-balancing regularizer that keeps prototype use spread out. On five standard
benchmarks at $\times2/\times3/\times4$, the resulting DPRNet matches or slightly
surpasses the strongest content-routing baseline at a comparable parameter and
computation budget, while the routing diagnostics give mechanism-level evidence that
the temperature stays unsaturated and the prototype usage becomes measurably more
balanced across all blocks. Because DPR preserves the external interface of the routing
block, it is not tied to this particular backbone and can serve as a generic,
confidence-aware content-routing module for other token-aggregation architectures.
Future work includes separating the contribution of each design element through
from-scratch retraining from a switch-neutral initialization and a multi-seed variance
study, completing the unified efficiency comparison under a single measurement protocol,
and extending DPR to larger training sets and real-world or arbitrary-scale degradations.
